% ═══════════════════════════════════════════════════════════════════
% EN Part 1: Introduction, principles, ladder, history, background
% ═══════════════════════════════════════════════════════════════════
\part{The Program and Its Principles}

\section{Introduction: the dynamic principle}\label{sec:intro}

This monograph unifies the results of a program developed step by
step through a series of computational stands: the calibration torus,
the K3 surface, the Klein quartic Jacobian, and the cyclotomic stand
$N=15/30$ on Fermat curves. Each rung of this ladder was originally
described in separate materials; here they are brought together,
made mutually consistent, and supplied with complete derivations, so
that the monograph reads as a single connected study --- from the
 founding principle to a closed system of lemmas and theorems with a
reproducible computational protocol. Every numerical statement is
certified: either by exact integer arithmetic or by direct
independent integration, reproducible with one command from the
\texttt{laboratory.py} lab.

The foundation of the program is the \emph{dynamic principle} ---
the rule by which the entire correction structure is built from
integer geometric data. The principle was fixed during the audit of
the original monographs and has never been violated since; its
formulation is simple.

\begin{keybox}{The dynamic principle (main formulation)}
\emph{Geometry supplies only the integer triple $(n,B,\lambda_0)$:
the order $n$ of a distinguished symmetry element, the integer
topological index $B$ of the carrier, and the spectral threshold
$\lambda_0$. All the remaining structure --- phases, braking terms,
effective corrections, normalizations, period lattices --- is built
from this triple by pure algebra and contains no free parameters.}
\end{keybox}

The principle is powerful precisely because of its thriftiness:
geometry does not ``guess'' or fit coefficients. It answers three
questions only: what order the rotation has, what integer index the
carrier has, and what the spectral threshold is. Everything else is
handled by a universal algebraic pipeline, identical on every rung
of the ladder. This is why the formulas on the torus, on K3, on the
Klein quartic, and on the Fermat curves are uniform: the formulas do
not repeat by analogy --- they are derived from one and the same
triple by one and the same rules.

A key methodological decision of the program is the rejection of
techniques that could create an appearance of results where there
are none. The final edition uses no PSLQ and no integer-relation
fitting for the claims of the main text; no Hecke character
L-values or externally postulated ``global phases''; no numerical
checks of what is proved identically; no appeals to self-adjoint
operators that never occur in the construction. Where numerical
verification is needed, it is performed by an independent method ---
direct quadrature over smooth pieces with no $\Ga$-functions at all
--- and serves as a certificate, not as a proof.

\subsection{What is proved and what is computed}

The program has accumulated a closed system of results, each proved
in the text of this monograph. For orientation, here is a summary;
full statements and proofs occupy the corresponding sections.

\begin{statusbox}[program summary]
\begin{itemize}[leftmargin=1.4em, itemsep=0.15em]
\item The Fermat curve genus $g=\tfrac{(N-1)(N-2)}{2}$ and the torus
  triple $(4,1,4\pi^2)$ --- proved from first principles
  (Sections~\ref{sec:torus},~\ref{sec:cycles}).
\item The N\'eron--Severi rank $\rho=20$ for the Fermat quartic
  $x^4+y^4+z^4+w^4=0$, signature $(1,19)$, sublattice determinant
  $64$ --- computed exactly over $\QQ$ from the explicit
  $49\times49$ intersection matrix (Section~\ref{sec:k3}).
\item Errata E8: on genus $3$, of the $64$ spin structures $36$ are
  even ($\mathrm{Arf}=0$) and $28$ odd ($\mathrm{Arf}=1$) --- proved by a complete Arf-invariant
  enumeration (Section~\ref{sec:e8}).
\item Certificates B, C, D, E, F, G, H --- accepted on all points of
  their protocols: closing the kernel $29$ by periods;
  $\mu_4$-equivariance $22=1+7+7+7$; the universal $\mu_4$ theorem
  $[L,P]=0$; flow termination with exact time $t^*=\lcm(\dots)$;
  rationalization with an a priori denominator bound $Q=256$; Hodge
  lattice indices and the Chowla--Selberg layer; the $N=15/30$ stand
  in the full Gross $\Ga$-normalization
  (Sections~\ref{sec:certB}--\ref{sec:klein}).
\item The closed period form $P_{a,b}(r,s)=\tfrac1N\zeta_N^{ra+sb}
  \Omega_{a,b}$, the Riemann relations as an identity, the
  $\Ga$-normalization with fractional shares of $\pi$, the
  polarization and computability of indices --- a closed system of
  lemmas and theorems of the $N=15/30$ stand
  (Sections~\ref{sec:cycles}--\ref{sec:indices}).
\item Protocol V1--V9: character census $91=1+6+84$ and
  $406=1+6+9+30+84+276$; closed-form deviations $3.9\cdot10^{-36}$;
  exact phases with deviation $0.0$; system rank $=g$ with condition
  numbers $8.6$ and $18.2$; deep record at $dps=70$ with deviation
  $1.7\cdot10^{-71}$ (Section~\ref{sec:protocol}).
\end{itemize}
\end{statusbox}

\subsection{How the monograph is organized}

The exposition is organized in large parts. Part~I (this section)
introduces the principle, the three corrections, and the ladder
methodology. Part~II describes the calibration stands: the torus,
where the whole mechanics is explicit; the K3 surface with maximal
N\'eron--Severi rank; the E8 errata on genus-3 spin structures; the
binary code as the bridge between geometry and cyclic algebra.
Part~III collects the program certificates --- from the cycle
certifier (certificate~A) to certificate~G on Hodge lattice indices;
the third stand, the Klein quartic Jacobian, is included here.
Part~IV is devoted to the $N=15/30$ stand: the most mature part of
the program, where the exposition is closed into a system of lemmas
and theorems relying on nothing external. Part~V contains the
V1--V9 protocol, summary tables, and appendices with complete
derivations.

A reader who wants to see the mechanics in action first should start
with Section~\ref{sec:torus} on the torus: the correction pipeline
takes a page and a half and every quantity is explicit. A reader
interested in the foundations should read sequentially: the cycle
combinatorics (Section~\ref{sec:cycles}) is built without any
analysis and serves as the base of everything else.

\section{The three corrections: origin and justification}
\label{sec:corrections}

The core of the framework consists of three corrections, each
derived rather than postulated. This section fixes their
definitions, justifies the origin of each, and assembles the united
formula. None of the quantities is a fitted parameter: all are
functions of the integer triple $(n,B,\lambda_0)$.

\subsection{The spin phase $\delta=\pi/N$}

The first correction is the spin phase. Suppose the geometry
supplies a symmetry element of order $N$: a rotation permuting the
sheets of the spin double cover. In the spin bundle, the lift of
this rotation has order $2N$ and permutes the two sheets; the phase
acquired by an eigenform upon going around the rotation axis is half
the elementary angle:
\begin{equation}\label{eq:delta}
  \delta \;=\; \frac{\pi}{N}.
\end{equation}
The origin of the formula is the double cover: on the base curve the
rotation has period $N$, but in the spin bundle a full circuit of
the axis returns the sheet only after two base circuits, hence the
elementary phase step is $\pi/N$, not $2\pi/N$. On the Klein curve
the order of the generator $C$ of $\Gamma(2,3,7)$ is $7$, giving the
phase $\delta_C=\pi/7$; on the torus the rotation order is $4$, so
$\delta=\pi/4$; on the Fermat curves the level plays the role of
$N$, and the phases are multiples of $\pi/15$ and $\pi/30$.

\begin{remark}[why the phase is not fitted]
Formula~\eqref{eq:delta} contains no free factors. Any deviation
would mean that the symmetry element has an order other than $N$;
but the order is an integer characteristic of the geometry. The
spin phase thus inherits the integrality of geometry: all further
corrections are built from $\delta=\pi/N$ algebraically.
\end{remark}

\subsection{Braking $\gamma=\delta^4/k$ and the effective phase
$\delta_{\mathrm{eff}}=\delta^5/k$}

The second correction is braking. The order-four phase correction
is suppressed by the carrier: the divisor is the integer
topological index $B$ --- the number $b_2$ of the carrier manifold
or another integer index of the rung. Definition:
\begin{equation}\label{eq:gamma}
  \gamma \;=\; \frac{\delta^4}{k},
  \qquad k = B\ \text{(carrier index)}.
\end{equation}
The universal choice of the framework is $k=b_2(K3)=22$: the second
Betti number of the K3 surface is an integer constant shared by all
rungs, since the K3 layer is present in the construction as the
carrier of corrections. On the torus the carrier is the torus itself
($k=1$), so $\gamma=\delta^4$; on the Klein curve $k=22$ is used; on
the $N=15/30$ stand there are several indices --- one per conductor
block --- and they are computed rather than postulated
(Section~\ref{sec:indices}).

The third correction combines the phase with braking into one
effective quantity:
\begin{equation}\label{eq:deff}
  \delta_{\mathrm{eff}} \;=\; \frac{\delta^5}{k} \;=\; \delta\cdot\gamma.
\end{equation}
For the Klein curve $(\pi/7)^5/22\approx1/1200$ --- the correction
is small, as required: the effective phase must be a perturbation,
not a replacement of the main term. On the torus with $k=1$ the
effective phase is large ($(\pi/4)^5\approx0.299$), and this is
exactly why the torus serves as the calibration stand: all effects
are visible at coarse precision and any construction error shows up
immediately.

\subsection{The united formula $\Delta_{Ch}$}

The three corrections are assembled into one quantity together with
the spectral threshold $\lambda_0$ and the rank term $R/4$:
\begin{keybox}{The united formula}
\begin{equation}\label{eq:deltaCh}
  \Delta_{Ch} \;=\; \lambda_0 \;-\; \frac{R}{4}
  \;+\; \frac{\delta^2}{2} \;-\; \frac{\delta^5}{k},
  \qquad
  b_{Ch}(n) \;=\; 1-\cos\frac{2\pi}{n}.
\end{equation}
Here $\lambda_0$ is the spectral threshold of the triple; $R$ is the
integer rank parameter of the rung; $\delta=\pi/n$ is the spin
phase; $k$ is the braking carrier index. The constant $b_{Ch}(n)$
is the radical characteristic of the rotation of order $n$, taking
exact radical values at the stand levels
(Theorem~\ref{th:normalization}).
\end{keybox}

Each term of formula~\eqref{eq:deltaCh} has its own origin. The term
$\lambda_0$ is the ``bare'' eigenvalue supplied by geometry; for the
Klein curve it is $\lambda_1(D^2_{\sigma_0})=3.338$, for the torus
$4\pi^2$. The term $-R/4$ is the rank term reflecting the size of
the unsplit part of the carrier. The term $+\delta^2/2$ is the main
second-order phase correction, universal for all rotations. The term
$-\delta^5/k$ is the fifth-order braking: small but mandatory
correction distinguishing the framework from a bare $\delta^2/2$.

Agreement with observed values is verified at every rung. For the
Klein curve the united formula gives $\Delta_{Ch}\approx3.4399$
against the observed $3.443$ --- a deviation of $0.0905\%$, fixed in
the protocol data (Table~\ref{tab:torus-triples}). The deviation is
not removed by fitting: it is a consequence of the same rules for
all rungs, and its smallness is evidence for the correctness of the
rules, not of a tuning.

\section{Ladder methodology}\label{sec:ladder}

The ladder of stands is a sequence of geometric objects of growing
complexity on which one and the same pipeline runs in full. Each
rung inherits the verified mechanisms of the previous one and adds
exactly one new phenomenon. This ordering guarantees that if
something is broken at a lower rung it is visible immediately, and
if a higher rung produces a surprise its source is localized by
contrast.

\begin{center}
\small
\begin{tabularx}{\textwidth}{@{}l l l >{\raggedright\arraybackslash}X@{}}
\toprule
Rung & $n$ & Genus & Constant and note \\
\midrule
Torus & $4$ & $1$ & triple $(4,1,4\pi^2)$; four spin structures;
  phase $\pi/4$; corrections derived from first principles \\
Klein curve & $7$ & $3$ & $b_{Ch}(7)\approx0.7526$; phases
  $\pi/2,\pi/3,\pi/7$; $B=b_2(K3)=22$; $j=-3375$ \\
Macbeath & $9$ & $7$ & $b_{Ch}(9)\approx0.1172$; Hurwitz tower \\
Hurwitz$_3$ & $11$ & $14$ & $b_{Ch}(11)\approx0.0823$; row aligned
  via $\mathrm{PSL}(2,13)$ \\
\textbf{Fermat (stand)} & $\mathbf{15}$ & $\mathbf{91}$ &
  $b_{Ch}(15)$ radical~\eqref{eq:b15}; phase lattice
  $\tfrac1{15}\ZZ[\zeta_{15}]$ \\
\textbf{Fermat (stand)} & $\mathbf{30}$ & $\mathbf{406}$ &
  $b_{Ch}(30)$ radical~\eqref{eq:b30}; phase lattice
  $\tfrac1{30}\ZZ[\zeta_{30}]$ \\
\bottomrule
\end{tabularx}
\end{center}

The ladder demonstrates a steady growth of scale: the genus grows
from $1$ to $406$, the phase lattices complicate from $\ZZ[i]$ to
$\ZZ[\zeta_{30}]$, and the number of integer indices grows --- one
per conductor block. The rules themselves never change: the spin
phase is always $\pi/n$, the braking always $\delta^4/k$, the
normalization always built from the periods themselves. The
constancy of the rules under scale growth is the main methodological
result of the ladder.

\begin{remark}[on the role of computational certificates]
Every quantitative statement of the monograph carries a certificate
of one of two types. Type one --- exact integer arithmetic: censuses,
ranks, determinants, indices are computed with no rounding and are
verified by an independent algorithm. Type two --- direct
independent integration: identities with $\Ga$-functions are checked
against tanh-sinh quadrature over smooth pieces in which
$\Ga$-functions never appear. The laboratory protocol reproduces all
certificates with one command; there are no ``peeked'' constants in
the code.
\end{remark}

\section{History of the program}\label{sec:history}

The program developed in several stages, each adding a new layer of
rigor. Understanding this evolution helps the reader see why the
final edition looks the way it does: every methodological decision
is dictated by specific control events recorded in the protocols.

\subsection{Stage I: the three corrections}

The original core was the set of three algebraic corrections
(Section~\ref{sec:corrections}): the spin phase $\delta=\pi/N$, the
braking $\gamma=\delta^4/k$, and the effective phase
$\delta_{\mathrm{eff}}=\delta^5/k$ with the united formula
\eqref{eq:deltaCh}. From the very start the integrality rule was in
force: every quantity must trace back to an integer characteristic
of the geometry. The rule was tested on the only rung available then
--- the Klein curve, where the triple $(7,22,\lambda_1)$ was taken
from the group $\Gamma(2,3,7)$ and the second coadjoint K3.

\subsection{Stage II: calibration stands}

The second stage added two calibration stands: the torus and K3. The
torus gave the complete explicit derivation of the corrections from
first principles (Section~\ref{sec:torus}) and revealed an important
fact: the $\mu_4$-equivariance for the symmetric encoding is
\emph{exact}, while for the $(7,22)$ encoding it is broken. This
contrast showed that the corrections inherit the geometry of the
carrier, and the naturality principle was formulated: a correction
must commute with the symmetry group. The principle was proved in
certificates~C and~D. The K3 stand (Section~\ref{sec:k3}) delivered
the first large exact computation: rank $20$, signature $(1,19)$,
determinant $64$ --- all over $\QQ$ with full control.

\subsection{Stage III: certificates}

The third stage introduced the institute of certificates: every
nontrivial step of the pipeline must have an independent check ---
either exact arithmetic or direct integration. The cycle certifier
(A) and certificate~B appeared, closing the dimension-$29$ kernel by
periods. Then followed C (naturality via $\mu_4$), D (the universal
$\mu_4$ theorem), E (flow termination with exact time), F
(rationalization with an a priori denominator bound), and G (Hodge
lattice indices, the Chowla--Selberg layer). The certificate
institute proved to be the most productive element of the program:
it turns every ``seems right'' into ``verified to $10^{-36}$'' or
``proved as an identity''.

\subsection{Stage IV: the cyclotomic stand and the closed system}

The fourth stage --- the $N=15/30$ stand (Part~V). Here the program
reached full maturity: the exposition is closed into a system of
lemmas and theorems relying neither on PSLQ, nor on L-values, nor on
numerical checks of the provable. All phases are derived explicitly
(Lemma~\ref{lem:closedform}), the Riemann relations are proved as an
identity (Theorem~\ref{th:riemann}), the normalization is reduced to
fractional shares of $\pi$ (Theorem~\ref{th:normalization}), and the
Hodge lattice indices are computed exactly
(Theorem~\ref{th:indexcomp}, certificate~H). Certificate~H completed
the arch: the full Gross $\Ga$-normalization of the volume gave
$\mathrm{vol}_h=1125=\sqrt{\disc}$ --- an integer, beautiful in its
arithmetic $3^4\cdot5^6$.

\subsection{Errata: two control events}

The history of the program includes two honestly recorded
corrections that themselves became results. First --- errata E8
(Section~\ref{sec:e8}): on genus $3$ there are $36$ even and $28$ odd spin structures,
not the swapped pair as printed in an early version; a full
enumeration of the $64$ structures by the Arf invariant closed the
question forever and added an automatic enumeration module to the
program. Second --- the E6 case on the torus: the correct derivation
gives $\Delta=0.7990$, not the previously printed $0.5964$; the
correction was obtained by the same chain of first principles and
entered all tables. Both events illustrated the power of the
protocol approach: any mismatch between the printed and the computed
is discovered automatically upon reproduction.

\section{Methodology: the five rules}\label{sec:rules}

Five rules have guided the program since stage~III. They are simple,
but they determine the architecture of both the monograph and the
laboratory.

\begin{keybox}{The five rules of the program}
\begin{enumerate}[leftmargin=2em, itemsep=0.2em]
\item \emph{The triple rule.} Geometry supplies only
  $(n,B,\lambda_0)$. Everything else is derived.
\item \emph{The closed-form rule.} Every checkable number is either
  derived in closed algebraic form or certified by direct
  independent measurement.
\item \emph{The certificate independence rule.} A certificate may
  not use the same algorithm as the computation it certifies;
  analysis is checked by integration without $\Ga$-functions,
  combinatorics by exact arithmetic.
\item \emph{The reproducibility rule.} Every result is reproducible
  by one command from a clean repository; randomness and seeds are
  forbidden.
\item \emph{The fixation rule.} Discrepancies are recorded as errata
  with full analysis; the fix enters the main text, not footnotes.
\end{enumerate}
\end{keybox}

Rule~1 separates geometry from algebra and makes the framework
portable: new geometry requires only the triple. Rule~2 eliminates
fitting: a number is either derived or measured independently.
Rule~3 protects against circularity: a certificate using the same
functions as the checked computation certifies nothing. Rule~4 makes
the monograph and the laboratory a single organism: every number of
the text has a source in \texttt{laboratory.py}. Rule~5 turns errors
into results: errata E8 and the E6 case are full sections with
automatic verification.

\subsection{Why computational certificates complement proofs}

One may ask: why certificates if the theorems are proved? The answer
is the division of labor between the analytic and the computational
channels. A theorem fixes an identity in general form; a certificate
guarantees that the concrete implementation --- code, libraries,
precision --- actually computes that identity and not a close
neighbor. A wrong sign, an index shift, a lost branch of a complex
logarithm --- classes of errors invisible in the statement but
immediately exposed by an independent numerical cross-check. On the
$N=15/30$ stand certificate~B does exactly this: the closed form
$\frac1N\zeta^{ra+sb}\Omega_{a,b}$ and the independent integration
agree to $10^{-36}$, hence the implementation is correct and the
phases are produced in the right order.

The symmetric situation occurs with integer structures: the
polarization type $(1,1,5,5,15,15,15,15)$ is obtained by the Smith
algorithm, but its arithmetic (the product of invariant factors
$=\disc$, the match with $3^4\cdot5^6$) is checked by a separate
integer line. Double registration --- algorithm plus arithmetic
invariant --- is the program's rule for all key tables.

\section{Mathematical background}\label{sec:background}

This section introduces all mathematical notions used by the
program, in a volume sufficient for independent reading of the
monograph. Each notion is defined, explained, and linked to its
concrete role in the pipeline. A reader familiar with Riemann
surfaces and Hodge structures may proceed directly to
Section~\ref{sec:torus}.

\subsection{Riemann surfaces and genus}

A \emph{Riemann surface} is a one-dimensional complex manifold: a
space locally modeled on a domain of the complex plane with
holomorphically compatible coordinate changes. Simplest examples ---
the Riemann sphere $\mathbb{P}^1$ (genus $0$), the torus
$\CC/\Lambda$ (genus $1$). The genus is an integer topological
invariant: the number of ``handles''; algebraically
$g=\dim H^0(C,\Omega^1)$ --- the dimension of the space of
holomorphic forms. For a smooth plane curve of degree $m$ the genus
is $\tfrac{(m-1)(m-2)}{2}$ --- this formula (Pl\"ucker) gives the
genus of the Klein quartic and, in another reading
(Lemma~\ref{lem:genus}), the genus of the Fermat curve.

The \emph{Fermat curve} $C_N=\{x^N+y^N=z^N\}$ is a plane curve of
degree $N$; its genus $(N-1)(N-2)/2$ grows quadratically: $91$ at
$N=15$ and $406$ at $N=30$. The curve carries the large symmetry
group $\mu_N\times\mu_N$: multiplication of coordinates by roots of
unity. This group organizes the entire computation --- both forms
and cycles decompose into its characters.

\subsection{Holomorphic forms and periods}

A \emph{holomorphic 1-form} is a differential $\omega=f(z)\,dz$ with
holomorphic $f$ in every chart. On a genus-$g$ curve their space is
$g$-dimensional; on the Fermat curve a basis is written explicitly
(Lemma~\ref{lem:forms}): $\eta_{i,j}=x^iy^j\,dx/y^{N-1}$.

A \emph{period} is the integral of a form over a closed cycle. Cycles
form the lattice $H_1(C,\ZZ)\cong\ZZ^{2g}$; a choice of cycle basis
turns the collection of periods into a \emph{period matrix} of size
$g\times2g$. The period matrix encodes the complex structure of the
curve recoverably (Torelli). On the Fermat curve every period has
the closed form $\frac1N\zeta^{ra+sb}\Omega_{a,b}$
(Lemma~\ref{lem:closedform}) --- the transcendental factor
$\Omega_{a,b}$ (a beta integral) and the algebraic phase (a root of
unity) are separated; this separation is the technical basis of the
whole $\Ga$-normalization.

The \emph{Jacobian} $\Jac(C)$ is the complex torus
$\CC^g/H_1(C,\ZZ)$ with its principal polarization; an abelian
variety into which the curve embeds via integration of forms (the
Abel--Jacobi map). For the Klein quartic the Jacobian is isogenous
to the cube of an elliptic curve with CM field $\QQ(\sqrt{-7})$ ---
Ribet's theorem, around which the stand of Section~\ref{sec:klein}
is built.

\subsection{Riemann relations and Hodge structures}

The period matrix in a symplectic basis must satisfy the \emph{Riemann
relations}: symmetry $\Omega^T=\Omega$ and positive definiteness of
$\im\Omega$. On a curve they are proved by Riemann's bilinear
identity (Theorem~\ref{th:riemann}) --- the only analytic tool the
program uses in this block. The space $H^1$ splits into the
\emph{Hodge structure} $H^1=H^{1,0}\oplus H^{0,1}$; the intersection
form defines the \emph{polarization}: the Riemann--Hodge relations
(Theorem~\ref{th:polarization}). The Hodge conjecture asserts that
rational classes orthogonal to all transcendental directions are
classes of algebraic cycles; the program builds the computable
infrastructure around this assertion --- polarizations, lattices,
indices --- and certifies each component exactly.

\subsection{Spin structures and the Arf invariant}

A \emph{spin structure} on a curve is a choice of square root of the
canonical bundle; on a genus-$g$ curve there are exactly $2^{2g}$ of
them. The parity of a structure is determined by the \emph{Arf
invariant} --- a quadratic form on $H^1(C,\ZZ/2)$. Parity controls
the theta constant: even structures have vanishing theta constant.
The even/odd counts by genus ($1/3$ at genus 1, $6/10$ at genus 2,
$28/36$ at genus 3) are classical; the program reproduces them by a
complete enumeration (Section~\ref{sec:e8}).

\subsection{Cyclotomic fields and roots of unity}

The \emph{roots of unity} $\zeta_N^k=e^{2\pi ik/N}$ form the cyclic
group $\mu_N$; the field $\QQ(\zeta_N)$ is the cyclotomic field of
degree $\phi(N)$. The phase lattice $\frac1N\ZZ[\zeta_N]$ is a
module over the ring of integers; the denominator $N$ is what
Theorem~\ref{th:normalization} fixes as ``fractional shares of
$\pi$''. The M\"obius function $\mu(d)$ underlies the census formula
\eqref{eq:mobius} --- a classical inclusion--exclusion scheme whose
exactness (integers!) makes the census a type-one certificate.

\subsection{Curves, lattices, and CM fields}

An elliptic curve with \emph{complex multiplication} (CM) has a
period ratio in an imaginary quadratic field; e.g.
$\tau=(1+\sqrt{-7})/2$ for the Klein Jacobian component. The lattice
$\ZZ\tau+\ZZ$ has Laplacian spectrum $|m\tau-n|^2$ --- the universal
formula of certificate~G1. The class numbers of small-discriminant
CM fields ($h(-7)=1$, $h(-3)=1$, $h(-15)=2$) make all mentioned
objects computable in integers: every CM object comes with its exact
minimal polynomial ($4X^2-7$, $4X^3-1$, etc.), and no quantity
remains ``just a number''.

\subsection{The tanh-sinh method: why it is independent}

The protocol verifies $\Ga$-forms by direct integration via tanh-sinh
(double-exponential quadrature). The method is applied to
\emph{smooth} integrands: the substitution $t=u^N/2$ on the two
halves of the interval removes the endpoint singularities of the
beta integral, and no $\Ga$-function ever appears in the quadrature
code --- this guarantees the independence of the certificate: an
error in the formula $\Ga(a/N)\Ga(b/N)/\Ga((a+b)/N)$ cannot support
itself, because the reference is computed by a completely different
algorithm. Agreement of the two methods to $10^{-36}$ at $dps=35$
certifies the correctness of the formula, not merely a numerical
coincidence.

\subsection{Smith normal form and indices}

The \emph{Smith normal form} (SNF) of an integer matrix is a
diagonal $d_1|d_2|\dots$ with exact integer invariant factors,
computed by integer elementary transformations with no rounding. For
polarization lattices the SNF gives the \emph{polarization type} ---
an integer invariant of the principally polarized variety. The
product of invariant factors equals the absolute determinant ---
for the $N=15/30$ stand this is
$1\cdot1\cdot5\cdot5\cdot15\cdot15\cdot15\cdot15=1265625=1125^2$,
and the agreement with the discriminant is an integer certificate
(Theorem~\ref{th:H1}). The sublattice index is the ratio of
determinants; the Hermite--Smith algorithm computes the saturation
basis exactly, making the block indices
(Theorem~\ref{th:indexcomp}) computable in the strict sense.
